%!TEX program = lualatex
% SPDX-License-Identifier: CC-BY-SA-4.0
% Requires: https://github.com/jemmybutton/fiziko (GPL-3.0)
% Set MPINPUTS to the directory containing fiziko.mp before running lualatex.
\documentclass[11pt]{article}

\usepackage{iftex}
\ifLuaTeX\else
  \PackageError{vintage-latex}{This example requires LuaLaTeX}
    {Run lualatex, not pdflatex or xelatex.}
\fi

\usepackage{fontspec}
\usepackage{unicode-math}
\usepackage{xcolor}
\usepackage{geometry}
\usepackage[tracking=true,protrusion=false,expansion=false]{microtype}
\usepackage{luamplib}

\setmainfont{EB Garamond}[Numbers={OldStyle,Proportional}]
\setmathfont{Garamond-Math.otf}
\definecolor{ragpaper}{HTML}{F3E9D6}
\definecolor{printerink}{HTML}{25201B}
\pagecolor{ragpaper}
\color{printerink}
\geometry{a4paper,top=22mm,bottom=22mm,left=24mm,right=24mm}
\pagestyle{empty}
\setlength{\parindent}{1.1em}

\mplibforcehmode
\everymplib{
  input fiziko.mp;
  randomseed := 2718;
  defineMinStrokeWidth(0.2pt);
  defineLightDirection(-1/6pi, 1/9pi);
  color inkcolor; inkcolor := (0.145, 0.125, 0.106);
  drawoptions(withcolor inkcolor);
  beginfig(1);
}
\everyendmplib{endfig;}

% One globe at a chosen station. fiziko carries its own coastline table and
% shades the sphere from the light direction declared above.
\newcommand{\station}[3]{% longitude, latitude, caption
  \begin{minipage}[t]{.31\linewidth}
    \centering
    \begin{mplibcode}
      draw globe(2.25cm, #1, #2);
      draw fullcircle scaled 4.5cm withpen pencircle scaled 0.55pt;
    \end{mplibcode}\par
    \vspace{.5em}
    {\small\itshape #3}
  \end{minipage}%
}

\begin{document}

\begin{center}
  {\small\textls[120]{\textsc{Plate I}}}\par
  \vspace{.3em}
  {\Large The Terrestrial Globe}\par
  \vspace{.2em}
  {\itshape shown from three stations, with the sun upon the left hand}
\end{center}

\vspace{1.5em}

\noindent
\station{40}{20}{Fig.\ 1. The Atlantic, between the Old World and the New.}
\hfill
\station{-100}{25}{Fig.\ 2. The Eastern hemisphere, from Africa to the Indies.}
\hfill
\station{0}{90}{Fig.\ 3. The northern regions, seen from above the pole.}

\vspace{1.6em}

\noindent\textls[70]{\textsc{The same body}} is here turned three times before
the observer. In each figure the illuminated part faces the left, whence the
light is supposed to come, and the shading grows heavier toward the limb on
the right, where the surface turns away both from the light and from the eye.
The outlines of the land are drawn upon the sphere and foreshortened with it,
so that a coast near the edge appears crowded together, while the same coast
brought to the centre is seen at its full extent.

\smallskip
The third figure is taken from a station above the pole. The meridians, which
in the other figures converge toward the top, here radiate from the centre
like the spokes of a wheel, and the lands of Europe, Asia, and America are
disposed around the frozen sea in a ring.

\vspace{1.2em}
\begin{center}
\begin{mplibcode}
  % A row of small globes turned by equal steps, as on the border of an
  % old map, showing how the same station passes across the disc.
  for i = 0 upto 5:
    draw globe(0.72cm, -30 + 60i, 20) shifted (2.05cm*i, 0);
  endfor;
  label.bot(btex \footnotesize\itshape The globe turned through sixty degrees
    between neighbours etex, (5.1cm, -1.1cm));
\end{mplibcode}
\end{center}

\vfill
\footnotesize\noindent
Every sphere on this page is one call of a fiziko macro drawn by LuaTeX's
built-in MetaPost. The captions, the lettering under the row, and the running
text are set together by LaTeX, so the plate and its explanation share one
font, one ink, and one paper.

\end{document}
